\section{Explicit Allocation Constructions}
\label{sec:allocation-methods}
We exploit the allocation freedom $z$ in two settings: a model can learn a new allocation from training data, or a frozen model can receive a native-relative redistribution. Both change adjacent log-frequency intervals. The first designs a density and samples its quantiles; the second distributes a prescribed extension across a transition band.
\subsection{An analytic construction from a tail-energy prior}
\label{sec:construction}

During learning, a model can adapt to a supplied allocation. The overlap analysis motivates spreading its finite rotary budget across scales while controlling concentration toward the slow end. We express this motivation as a geometry-inspired, weight-independent density-design prior, using a density over normalized exponent space $[0,1]$. Let $S_\rho(t)=\int_t^1\rho(\phi)\,d\phi$, and consider the convex design criterion
\begin{equation}
\Capp[\rho]=\frac\alpha2\int_0^1\rho(\phi)^2\,d\phi+\frac\beta2\int_0^1S_\rho(t)^2\,dt,
\qquad \text{s.t. } \int_0^1\rho=1,\ \rho\ge0.
\label{eq:Capp}
\end{equation}
The first term penalizes concentration; the second penalizes cumulative slow-tail mass. Their balance selects how density shifts toward smaller exponents. The tail term is the quadratic form of the Dirichlet--Neumann Green kernel $\min(\phi,\psi)$ (Appendix~\ref{sec:proofs}).

For $\alpha>0$, $\beta\ge0$, this criterion has a unique positive minimizer, $\rho_\tau(\phi)=\tau\cosh[\tau(1-\phi)]/\sinh\tau$, where $\tau=\sqrt{\beta/\alpha}$. Its inverse CDF gives
\begin{equation}
Q_\tau(u)=1-\frac{\operatorname{arcsinh}((1-u)\sinh\tau)}{\tau},\qquad u_k=\frac{k+1/2}{K}.
\label{eq:warp}
\end{equation}
The uniform-density limit is recovered as $\tau\to0$. Theorem~\ref{thm:ode} in Appendix~\ref{sec:proofs} gives the complete construction proof.

\evq{} converts the optimal density into a finite table through \eqref{eq:warp}. We use Cosh for the analytic shape and state the grid and endpoint treatment for each finite EVQ-Cosh table. For $\tau>0$, $\rho_\tau$ is decreasing, so its inverse CDF $Q_\tau$ is convex. For equally spaced nodes $u_k$, affine endpoint anchoring therefore gives
\begin{equation}
z_{\tau,k}=\frac{Q_\tau(u_k)-Q_\tau(u_0)}{Q_\tau(u_{K-1})-Q_\tau(u_0)}
\le \frac{k}{K-1},\qquad
\omega_{\tau,k}\ge\omega_{\mathrm{Geo},k}
\label{eq:anchored-cosh-order}
\end{equation}
at fixed $(a,R)$: convexity places each quantile below the chord between the extreme nodes. Interior pairs move toward faster frequencies while endpoints remain fixed (Fig.~\ref{fig:spectral-budget-overview}a). Once $\tau$, grid, and endpoint anchoring are specified, installation evaluates the table once without changing the rotary operator or trainable parameter count.

The strength $\tau$ sets this trade-off. We use declared reference settings and compare neighboring strengths and a second analytic family in Appendix~\ref{sec:exact-range-factorial}; Appendix~\ref{sec:tau-scaling} gives the operating rules and their empirical basis.
